\chapter{Reproducibility of the Theoretical Manuscript}

This appendix records how to rebuild the thesis artifact. It intentionally does
not define a simulation protocol. The present manuscript is the theoretical
part of the dissertation.

\section{Source Layout}

The LaTeX source lives under:

\begin{verbatim}
reports/thesis/
\end{verbatim}

The main file is:

\begin{verbatim}
reports/thesis/main.tex
\end{verbatim}

The chapter files are:

\begin{verbatim}
reports/thesis/chapters/
\end{verbatim}

The bibliography is:

\begin{verbatim}
reports/thesis/references.bib
\end{verbatim}

\section{Compile the Thesis}

From the repository root, the manual build is:

\begin{verbatim}
cd reports/thesis
pdflatex -interaction=nonstopmode -halt-on-error main.tex
bibtex main
pdflatex -interaction=nonstopmode -halt-on-error main.tex
pdflatex -interaction=nonstopmode -halt-on-error main.tex
\end{verbatim}

The final PDF is copied to:

\begin{verbatim}
outputs/pdf/thesis_draft.pdf
\end{verbatim}

\section{Evidence Boundary}

This theoretical manuscript does not use plots or time-domain traces as
evidence. Numerical studies, when added later, should be treated as a separate
validation layer. A future validation artifact should record:

\begin{enumerate}
\item the full dynamic model and operating point;
\item the retained and condensed state partition;
\item the construction of \(\Pi(s)\), its channel decomposition, and any
\(\Sigma(s)\) approximation;
\item full-model eigenvalue targets;
\item accepted error tolerances;
\item failure cases and rejected hypotheses.
\end{enumerate}

This boundary protects the thesis from mixing mathematical derivation with
uncalibrated simulation evidence.
